%!TEX root = ../../thesis-trier.tex

Here, in \cref{sync_max:thm:abJordan},
we characterize $\{a,b\}$-primitivity 
in terms of generating sets for the full symmetric group, thereby generalizing Jordan's theorem (\cref{sync_max:thm:prim_char:jordan}).
As a consequence, in \cref{sync_max:cor:cycletransgenerateSn},
we characterize those pairs consisting of a permutation
and a transposition that generate the full symmetric group.
Clearly, adding permutations preserves $A$-primitivity.
\Cref{sync_max:lem:notSn} gives a condition for the preservation of $A$-primitivity when removing permutations.

\begin{lemma}
\label{sync_max:lem:notSn}	
	Let $G \le_G \mathcal S_n$ and $t \in \mathcal S_n$ be a permutation. 
	If $\langle G \cup \{t\} \rangle_G$
	is $A$-primitive, then $G$ is also $A$-primitive
	for $A = \move{t}$.
\end{lemma}
\begin{proof}
	Suppose $G$ is not $A$-primitive.
	Then there exists a non-trivial block $\Delta \subseteq [n]$
	such that $A \subseteq \Delta$.
	Let $x \in A$. Then $t(x) \ne x$.
	Since~$t$ is a permutation, we find $t(t(x)) \ne t(x)$
	and so $t(x) \in A$. Thus $t(A) \subseteq A$.
	We also have $t(t^{-1}(x)) = x \ne t^{-1}(x)$,
	giving $t^{-1}(x) \in A$
	and then $x \in t(A)$. Thus $A \subseteq t(A)$.
	As all points outside of $A$ are fixed by $t$,
	we find  $\applyfunc{\Delta}{t} = \Delta$.
	This yields that $\Delta$
	is also a block for $\langle G \cup \{t\} \rangle_G$
	and so this group is not $A$-primitive.
\end{proof}

By \cref{sync_max:lem:notSn}, if $g$ and $h$ generate a primitive group,
then $\langle g \rangle_G$ is $\move{h}$-primitive.

\begin{theorem}\label{sync_max:thm:abJordan}
	Let $a,b \in [n]$ be distinct.
	Then $G \le_G \mathcal S_n$ is $\{a,b\}$-primitive
	if and only if $\langle G \cup \{ (a ~ b) \} \rangle_G = \mathcal S_n$.
\end{theorem}
\begin{proof}
	First, suppose $\langle G \cup \{ (a ~ b) \} \rangle_G = \mathcal S_n$.
	As the full symmetric group $\mathcal S_n$ is $\{a,b\}$-primitive,
	by \cref{sync_max:lem:notSn} we find that $G$ is $\{a,b\}$-primitive.		
	
	Conversely, suppose $G$ is $\{a,b\}$-primitive.
	Set $T =  \{ (\applyfunc{a}{g} ~ \applyfunc{b}{g}) \mid g \in G \}$.
	The graph with vertex set $[n]$ and edge set
	$
	\{ \{c,d\} \mid (c ~ d) \in T\}
	$
	equals the undirected orbital graph on $\{a,b\}$.
	By \cref{syncmax:prop:abprim_undirectedorbital} it is connected.
	So, by \cref{prop:connected_transpositions_gen_Sn},
	we get $\langle T \rangle_G = \mathcal S_n$.
	By \cref{prelim:conj}, we have $(\applyfunc{a}{g} ~ \applyfunc{b}{g}) = (a ~ b)^g$,
	which gives $T \subseteq \langle G \cup \{ (a ~ b) \} \rangle_G$.
	Hence $\mathcal S_n = \langle G \cup \{ (a ~ b) \} \rangle_G$.		
\end{proof}

The fact that $(0 ~ 1)$ and $(0 ~ 1 ~ \ldots ~ n - 1)$
generate $\mathcal S_n$ is well-known, for example, 
see~\cite{cameron_1999,dixon1996permutation}.
As a consequence of \cref{sync_max:thm:abJordan,syncmax:prop:cyclic-group-abprim},
we obtain the following characterization
of those permutations that, together with
a transposition, generate the whole symmetric group.

\begin{corollary}\label{sync_max:cor:cycletransgenerateSn}
	Let $g \in \mathcal S_n$ and $a,b \in [n]$ be distinct.
	Then $\langle \{ g, (a ~ b)\} \rangle_G = \mathcal S_n$
	if and only if $g$ induces two cycles of coprime length, one cycle moving~$a$ and the other~$b$,
	or $g$ cyclically permutes
	its domain and  $g^i(a) = b$
	for some $i \ge 0$ coprime to $n$. %for some (and hence all) $i > 0$ with
\end{corollary}